\section[Finite Native Bn Fingerprint Boundary]{Finite Native $B_n$ Fingerprint Boundary}\label{app:fullfingerprint}

The main theorem of this paper is not a full native $B_n$ Fourier fingerprint theorem.  We nevertheless record one
finite boundary computation because it is useful for interpretation: the naive expectation that every native irreducible
block saturates the $R=+$ bound already fails at $n=8$.

All statements in this appendix section are finite statements for the displayed $n=8$ atlas.  They do not imply an
all-$n$ classification, an all-irrep saturation theorem, or a general Johnson-scheme vanishing theorem.  The external
$n=8,k\ge2$ defect census contains $41$ fixed-plus family-instances and $114$ defect rows.  The currently certified
accounting covers $6$ family-instances and $44$ rows, leaving $35$ family-instances and $70$ rows open.

\begin{center}
\begingroup
\scriptsize
\setlength{\tabcolsep}{4pt}
\begin{tabular}{c|r|r|r|r|r|r}
\textbf{Slice} & \textbf{fam.} & \textbf{rows} & \textbf{acct. fam.} & \textbf{acct. rows} & \textbf{open fam.} & \textbf{open rows}\\
\hline
$n=8,k=2$ & 5 & 20 & 1 & 15 & 4 & 5\\
$n=8,k=3$ & 36 & 94 & 5 & 29 & 31 & 65\\
\hline
total & 41 & 114 & 6 & 44 & 35 & 70
\end{tabular}
\endgroup
\end{center}

The accounted family-instances are
\begin{align*}
(n,k,\lambda^+)\in\{&(8,2,(1)),\ (8,3,(1,1)),\ (8,3,(2)),\\
&(8,3,(4)),\ (8,3,(4,1)),\ (8,3,(2,2,2,1,1))\}.
\end{align*}
The phrase ``accounted'' here means finite exact certificate or finite quotient-layer accounting, not an all-$n$ theorem.

\begin{proposition}[Pre-Specht cancellation for the $(1,1)$ zero family]\label{prop:lambda11zero}
For $X=\BtildeI_8(3,3)$ and every partition $\nu\vdash6$,
\[
  M_{(1,1)\mid\nu}(X)=0.
\]
\end{proposition}

\begin{proof}[Certified proof sketch]
Before applying a complement Specht module, the fixed-plus block with $\lambda^+=(1,1)$ is a matrix over the group
algebra $\mathbb Q[S_6]$.  The exact certificate groups the $192$ active signed terms into $96$ group-algebra keys.
For every key there are two surviving completions with the same complement permutation and opposite $S_2$ sign; their
coefficients are $-4$ and $4$.  Thus all $96$ keys vanish already in $\mathbb Q[S_6]$.  Applying any representation
$S^\nu$, $\nu\vdash6$, preserves zero.
\end{proof}

\begin{proposition}[Pure-plus Johnson degree-drop]\label{prop:purepluszero}
For $X=\BtildeI_8(3,3)$,
\[
  M_{(2,2,2,1,1)\mid\varnothing}(X)=0.
\]
\end{proposition}

\begin{proof}[Certified proof sketch]
The conjugate partition of $(2,2,2,1,1)$ is $(5,3)$, so
\[
  S^{(2,2,2,1,1)} \cong S^{(5,3)}\otimes\operatorname{sgn}.
\]
After the corresponding sign twist, the certified operator $A_{\operatorname{sgn}}$ acts on the $3$-subset permutation
module $M^{(3)}$.  Let $U_2$ be the pair-incidence submodule spanned by the vectors $v_P$ of $3$-subsets containing a
fixed pair $P$.  The Johnson certificate verifies the exact containment
\[
  A_{\operatorname{sgn}}(M^{(3)})\subseteq U_2.
\]
Since the top quotient $M^{(3)}/U_2$ is isomorphic to $S^{(5,3)}$, the sign-twisted operator kills this quotient.
Untwisting by the sign representation gives the displayed pure-plus zero.  An independent direct rational seminormal
Specht computation also gives zero entries throughout the $28$-dimensional block.
\end{proof}

\begin{remark}[Appendix boundary]
The finite zero mechanisms above explain all twelve zero rows under the $n=8,k=3$ atlas census: eleven rows from
Proposition~\ref{prop:lambda11zero} and one row from Proposition~\ref{prop:purepluszero}.  Other certified finite
packages account for selected nonzero extra-defect families, including $T_1^{(8,2)}$, $T_2^{(8,3)}$, $T_4^{(8,3)}$, and
$\lambda^+=(4,1)$ at $n=8,k=3$.  The remaining extra-defect rows are not classified here.
\end{remark}
